\documentclass[10pt,a4paper]{article}

\usepackage[margin=14mm,top=14mm,bottom=16mm]{geometry}
\usepackage{fontspec}
\usepackage[table]{xcolor}
\usepackage{tabularx}
\usepackage{array}
\usepackage{enumitem}
\usepackage{titlesec}
\usepackage{fancyhdr}
\usepackage{lastpage}
\usepackage{ragged2e}
\usepackage{amssymb}
\usepackage[most]{tcolorbox}
\usepackage[hidelinks]{hyperref}

\IfFontExistsTF{Arial}{
  \setmainfont{Arial}
  \setsansfont{Arial}
}{
  \setmainfont{TeX Gyre Heros}
  \setsansfont{TeX Gyre Heros}
}

\definecolor{AIEPRed}{HTML}{D62027}
\definecolor{AIEPNavy}{HTML}{082B5C}
\definecolor{AIEPInk}{HTML}{243B53}
\definecolor{AIEPSlate}{HTML}{52606D}
\definecolor{AIEPBlue}{HTML}{1B6CA8}
\definecolor{AIEPCyan}{HTML}{35B7C6}
\definecolor{AIEPGreen}{HTML}{2E8B57}
\definecolor{AIEPGold}{HTML}{D8A928}
\definecolor{AIEPPaper}{HTML}{F8F3EC}
\definecolor{AIEPSoftBlue}{HTML}{E8F1F8}
\definecolor{AIEPGreenSoft}{HTML}{EAF5EF}
\definecolor{AIEPGoldSoft}{HTML}{FFF5D8}
\definecolor{Line}{HTML}{D2DAE3}
\definecolor{White}{HTML}{FFFFFF}

\pagestyle{fancy}
\fancyhf{}
\renewcommand{\headrulewidth}{0pt}
\fancyfoot[L]{\scriptsize\textcolor{AIEPSlate}{Planificación docente - Taller 3}}
\fancyfoot[R]{\scriptsize\textcolor{AIEPSlate}{GeoGreen Escolar Osorno - \thepage/\pageref{LastPage}}}

\setlength{\parindent}{0pt}
\setlength{\parskip}{3.7pt}
\setlist[itemize]{leftmargin=4.3mm,topsep=1pt,itemsep=1.2pt,parsep=0pt}
\setlist[enumerate]{leftmargin=5mm,topsep=1pt,itemsep=1.2pt,parsep=0pt}
\renewcommand{\arraystretch}{1.18}
\titleformat{\section}{\Large\bfseries\color{AIEPNavy}}{}{0pt}{}
\titleformat{\subsection}{\normalsize\bfseries\color{AIEPInk}}{}{0pt}{}

\newcolumntype{Y}{>{\RaggedRight\arraybackslash}X}
\newcolumntype{L}[1]{>{\RaggedRight\arraybackslash}p{#1}}

\newtcolorbox{card}[1][]{
  enhanced,colback=White,colframe=Line,boxrule=0.5pt,arc=1.2mm,
  left=3.2mm,right=3.2mm,top=2.8mm,bottom=2.8mm,#1
}
\newtcolorbox{softcard}[2][]{
  enhanced,colback=#2!10,colframe=#2,boxrule=0.55pt,arc=1.2mm,
  left=3.2mm,right=3.2mm,top=2.8mm,bottom=2.8mm,#1
}
\newcommand{\kicker}[1]{%
  {\color{AIEPRed}\scriptsize\bfseries\MakeUppercase{#1}}\par\vspace{2.5mm}
  {\color{AIEPRed}\rule{13mm}{1.2pt}}\vspace{3mm}
}
\newcommand{\labeltext}[1]{%
  {\scriptsize\bfseries\color{AIEPSlate}\MakeUppercase{#1}}\par\vspace{1.2mm}
}
\newcommand{\minute}[1]{%
  \fcolorbox{AIEPNavy}{AIEPNavy}{\textcolor{White}{\bfseries\strut #1}}
}

\begin{document}

% 1 ---------------------------------------------------------------------------
\thispagestyle{fancy}
\kicker{GeoGreen Escolar Osorno}
{\Huge\bfseries\color{AIEPNavy} Taller 3 - De la idea al prototipo}\par
\vspace{2mm}
{\large\color{AIEPSlate} Planificación lectiva para facilitar una experiencia de sensores, prototipado y desarrollo asistido por inteligencia artificial.}\par
\vspace{6mm}

\begin{tabularx}{\textwidth}{Y Y Y}
\begin{softcard}[height=33mm,valign=top]{AIEPBlue}
\labeltext{Duración}
{\Large\bfseries\color{AIEPNavy}90 minutos}\par
Sesión presencial por bloque.
\end{softcard}
&
\begin{softcard}[height=33mm,valign=top]{AIEPCyan}
\labeltext{Organización}
{\Large\bfseries\color{AIEPNavy}5 equipos}\par
6 estudiantes por equipo.
\end{softcard}
&
\begin{softcard}[height=33mm,valign=top]{AIEPGreen}
\labeltext{Producto}
{\Large\bfseries\color{AIEPNavy}1 propuesta}\par
Tecnológica inicial por equipo.
\end{softcard}
\end{tabularx}

\vspace{4mm}
\begin{card}
\labeltext{Datos de ejecución}
\begin{tabularx}{\textwidth}{@{}L{36mm}Y@{}}
\textbf{Fecha} & Lunes 24 de agosto de 2026.\\
\textbf{Bloque A} & 09:00 a 10:30.\\
\textbf{Bloque B} & 10:45 a 12:15.\\
\textbf{Participantes} & 60 estudiantes de tercero medio, distribuidos en dos bloques de 30.\\
\textbf{Responsabilidad} & Escuela de Ingeniería, Energía y Tecnología (I/E/T).\\
\textbf{Continuidad} & Cuatro mentorías de desarrollo y evento final del programa.\\
\end{tabularx}
\end{card}

\begin{softcard}{AIEPRed}
\labeltext{Sentido formativo}
GeoGreen se presenta como un caso real de innovación: una necesidad observable se transforma en
datos, reglas, alertas, software y un producto físico. El taller no busca convertir la sesión en una
clase extensa de electrónica; busca que cada equipo comprenda el ciclo completo, explore tecnología
de forma segura y formule una idea propia que pueda desarrollar después.
\end{softcard}

\begin{softcard}{AIEPGold}
\labeltext{Pregunta movilizadora}
{\large\bfseries\color{AIEPNavy}¿Qué situación de nuestro entorno podríamos observar con un sensor para crear una respuesta útil?}
\end{softcard}

\begin{card}
\labeltext{Principio de facilitación}
Mostrar posibilidades abre el horizonte; las preguntas, decisiones y pruebas pertenecen a los
equipos. La docencia orienta, hace visible el razonamiento técnico y resguarda la seguridad, sin
resolver el proyecto por ellos.
\end{card}

\newpage
% 2 ---------------------------------------------------------------------------
\kicker{Propósito y resultados}
{\LARGE\bfseries\color{AIEPNavy}Qué deben comprender, decidir y producir}\par
\vspace{4mm}

\begin{softcard}{AIEPBlue}
\labeltext{Objetivo general}
Comprender cómo una idea puede convertirse en un prototipo tecnológico mediante la relación entre
problema, variable observable, sensor, programación, evidencia y evolución del producto.
\end{softcard}

\begin{tabularx}{\textwidth}{Y Y}
\begin{card}[height=69mm,valign=top]
\labeltext{Objetivos específicos}
Al finalizar, las y los estudiantes podrán:
\begin{itemize}
  \item reconocer el ciclo \textbf{sensar → procesar → comunicar → actuar};
  \item diferenciar placa, sensor, actuador, protoboard y software;
  \item seleccionar un sensor coherente con una variable observable;
  \item formular una regla inicial en formato \textbf{“cuando..., entonces...”};
  \item usar IA como apoyo para explorar, explicar y prototipar, verificando sus respuestas;
  \item definir una prueba segura y un siguiente paso verificable.
\end{itemize}
\end{card}
&
\begin{card}[height=69mm,valign=top]
\labeltext{Aprendizajes clave}
\begin{itemize}
  \item Un sensor no “resuelve” el problema: entrega una señal que debe interpretarse.
  \item Un prototipo sirve para reducir incertidumbre mediante pruebas.
  \item Simular permite corregir antes de conectar hardware real.
  \item La IA acelera el proceso, pero no sustituye la comprobación humana.
  \item La calidad crece por iteraciones: idea, circuito, código, evidencia, interfaz y producto.
\end{itemize}
\end{card}
\end{tabularx}

\begin{softcard}{AIEPGreen}
\labeltext{Salida obligatoria por equipo}
Una \textbf{ficha de propuesta tecnológica inicial} que contenga: problema y contexto; personas
afectadas; variable observable; sensor elegido y fundamento; regla de funcionamiento; primera
evidencia o prueba pendiente; nivel de madurez; roles; y próximo hito verificable.
\end{softcard}

\begin{card}
\labeltext{Entradas que el equipo recupera}
\begin{tabularx}{\textwidth}{@{}L{39mm}Y@{}}
\textbf{Desde Taller 1} & Problema ambiental priorizado y contexto en que ocurre.\\
\textbf{Desde Taller 2} & Material o residuo analizado y hallazgos de su ficha técnica.\\
\textbf{En Taller 3} & Una idea tecnológica que conecte la necesidad con una variable medible.\\
\end{tabularx}
\end{card}

\begin{softcard}{AIEPGold}
\labeltext{Indicador de logro}
El equipo puede explicar, en 30 a 40 segundos, \textbf{qué quiere observar, con qué sensor, qué hará
el sistema con ese dato y cuál será su siguiente prueba}.
\end{softcard}

\newpage
% 3 ---------------------------------------------------------------------------
\kicker{Metodología y condiciones}
{\LARGE\bfseries\color{AIEPNavy}Explorar con autonomía, trabajar con seguridad}\par
\vspace{4mm}

\begin{tabularx}{\textwidth}{Y Y}
\begin{card}[height=73mm,valign=top]
\labeltext{Enfoque metodológico}
\begin{itemize}
  \item \textbf{Caso inspirador:} GeoGreen muestra una trayectoria posible, no una respuesta que deba copiarse.
  \item \textbf{Demostración breve:} cada concepto aparece conectado a una función visible.
  \item \textbf{Aprender haciendo:} los equipos manipulan, preguntan, comparan y documentan.
  \item \textbf{Desarrollo asistido:} la IA ayuda a investigar, proponer y depurar.
  \item \textbf{Verificación humana:} pinout, voltaje, conexiones, código y resultados se comprueban.
  \item \textbf{Progresión:} el taller inicia la propuesta; el desarrollo continúa entre mentorías.
\end{itemize}
\end{card}
&
\begin{card}[height=73mm,valign=top]
\labeltext{Rol de quien facilita}
\begin{itemize}
  \item contextualizar el proyecto y hacer visible el ciclo completo;
  \item traducir conceptos técnicos con ejemplos concretos;
  \item modelar preguntas útiles para trabajar con IA;
  \item revisar conexiones antes de energizar;
  \item acompañar decisiones sin apropiarse de la idea;
  \item proteger el tiempo de trabajo de los equipos;
  \item recoger evidencia y asegurar la salida obligatoria.
\end{itemize}
\end{card}
\end{tabularx}

\begin{softcard}{AIEPRed}
\labeltext{Protocolo de seguridad antes del USB}
\begin{enumerate}
  \item Mantener la placa \textbf{sin alimentación} mientras se arma o modifica el circuito.
  \item Confirmar el \textbf{pinout}, la polaridad, el voltaje y la tierra común.
  \item Usar la \textbf{resistencia correspondiente} con cada LED.
  \item Solicitar revisión docente antes de conectar el cable USB.
  \item Si la placa o un componente se calienta, desconectar de inmediato y avisar.
\end{enumerate}
\textbf{Precisión técnica:} el HC-SR04 entrega 5 V en Echo. En UNO R4 WiFi puede conectarse de forma
directa; en una placa ESP32 de 3,3 V necesita divisor de tensión o adaptación de nivel.
\end{softcard}

\begin{card}
\labeltext{Materiales a disponer por bloque}
\begin{itemize}
  \item computador y proyector; presentación, videos e infografías del taller;
  \item prototipo GeoGreen y, si es posible, acceso al dashboard para mostrar el ciclo completo;
  \item cinco estaciones de exploración con placa, protoboard, cable USB, jumpers y sensores del banco;
  \item componentes de salida disponibles: LEDs, resistencias, buzzer o pantalla;
  \item cinco copias de la ficha de propuesta tecnológica inicial y lápices;
  \item equipos con acceso a Wokwi y a una herramienta de IA cuando la conectividad lo permita.
\end{itemize}
\end{card}

\begin{softcard}{AIEPCyan}
\labeltext{Organización del equipo}
Se mantienen seis responsabilidades: coordinación, problema y propósito, investigación, prototipo
físico, programación y datos, y relato/evidencia. Las responsabilidades ordenan el trabajo; las
decisiones relevantes se toman como equipo.
\end{softcard}

\newpage
% 4 ---------------------------------------------------------------------------
\kicker{Secuencia de 90 minutos}
{\LARGE\bfseries\color{AIEPNavy}Mapa operativo de la sesión}\par
\vspace{4mm}

\begin{tabularx}{\textwidth}{@{}L{22mm}L{48mm}Y@{}}
\rowcolor{AIEPNavy}
\textcolor{White}{\bfseries Tiempo} &
\textcolor{White}{\bfseries Momento} &
\textcolor{White}{\bfseries Acción y evidencia}\\
\minute{0–5} &
\textbf{Apertura} &
Presentar el desafío: observar una situación real y convertirla en una respuesta tecnológica útil. Recuperar los equipos y sus productos anteriores.\\[1.8mm]
\minute{5–15} &
\textbf{Bloque 1}\newline Una idea puede crecer &
Mostrar GeoGreen como proyecto emblemático y recorrer su evolución: necesidad, sensor, semáforo y buzzer, WiFi, visualización, carcasa y PCB. Cerrar con la pregunta movilizadora.\\[1.8mm]
\minute{15–25} &
\textbf{Bloque 2}\newline Del mundo físico al dato &
Explicar placa, protoboard, sensor, actuador y flujo entrada–proceso–salida. Hacer visible cómo se conectan los rieles y las filas de la protoboard.\\[1.8mm]
\minute{25–35} &
\textbf{Demostración segura} &
Modelar el protocolo previo al USB y una lectura simple. Comparar UNO R4 WiFi y ESP32 solo desde sus casos de uso y voltajes relevantes.\\[1.8mm]
\minute{35–45} &
\textbf{Bloque 3}\newline Elegir qué observar &
Cada equipo recupera su problema, identifica una variable y selecciona un sensor. Si no existe una elección clara, explora el banco disponible y justifica una alternativa.\\[1.8mm]
\minute{45–60} &
\textbf{Prototipar con IA} &
Formular una consulta con contexto, placa, sensor, objetivo y restricciones. Revisar la propuesta; simular o preparar una conexión física segura. Registrar decisiones y correcciones.\\[1.8mm]
\minute{60–70} &
\textbf{Primera evidencia} &
Obtener una lectura, una simulación, un esquema revisado o dejar descrito el próximo ensayo exacto. Completar la regla “cuando..., entonces...”.\\[1.8mm]
\minute{70–80} &
\textbf{Bloque 4}\newline Del prototipo al producto &
Presentar la escalera de madurez: circuito, código, comunicación, software, carcasa y PCB. Conectar cada nivel con una pregunta de mejora.\\[1.8mm]
\minute{80–87} &
\textbf{Propuesta y siguiente hito} &
Completar la ficha, verificar que las seis responsabilidades estén cubiertas y acordar un avance realizable antes de la Mentoría 1.\\[1.8mm]
\minute{87–90} &
\textbf{Cierre relámpago} &
Dos o tres equipos comparten su propuesta en 30–40 segundos. Recoger todas las fichas y cerrar con la ruta de continuidad.\\
\end{tabularx}

\vspace{4mm}
\begin{softcard}{AIEPGold}
\labeltext{Regla de manejo del tiempo}
La exposición habilita la acción. Si un momento se extiende, recortar explicación secundaria antes
que el tramo de trabajo de equipos. La sesión debe terminar con una decisión y un siguiente paso,
no solo con información.
\end{softcard}

\newpage
% 5 ---------------------------------------------------------------------------
\kicker{Mediación por bloque}
{\LARGE\bfseries\color{AIEPNavy}Preguntas que hacen avanzar el razonamiento}\par
\vspace{4mm}

\begin{softcard}{AIEPBlue}
\labeltext{Bloque 1 · GeoGreen como caso de innovación}
\textbf{Preguntas:} ¿qué problema observa?, ¿qué fue necesario agregar para que el dato se entendiera?,
¿qué parte pertenece al hardware y cuál al software?, ¿qué mejora aporta valor y cuál solo decora?\par
\textbf{Señal de progreso:} el grupo identifica que un proyecto tecnológico puede evolucionar por
capas y que su propia idea no necesita comenzar “terminada”.
\end{softcard}

\begin{softcard}{AIEPCyan}
\labeltext{Bloque 2 · Del mundo físico al dato}
\textbf{Preguntas:} ¿qué variable cambia?, ¿qué componente la detecta?, ¿qué señal recibe la placa?,
¿qué salida permitiría entender el resultado?, ¿qué debemos comprobar antes de energizar?\par
\textbf{Señal de progreso:} el grupo puede narrar el flujo entrada → decisión → salida y reconoce
qué conexiones están realmente unidas en la protoboard.
\end{softcard}

\begin{softcard}{AIEPGreen}
\labeltext{Bloque 3 · Prototipar con sensores y agentes}
\textbf{Preguntas:} ¿qué quieren medir exactamente?, ¿por qué este sensor y no otro?, ¿qué placa
están usando?, ¿qué dato falta para confiar en la respuesta de la IA?, ¿cómo sabrán que la prueba funcionó?\par
\textbf{Señal de progreso:} el equipo justifica sensor y regla, realiza una verificación técnica y
obtiene evidencia o deja definido un ensayo concreto.
\end{softcard}

\begin{softcard}{AIEPGold}
\labeltext{Bloque 4 · Del prototipo al producto}
\textbf{Preguntas:} ¿en qué nivel de madurez está la idea?, ¿qué debería ver una persona usuaria?,
¿cómo viajaría el dato?, ¿qué protegería el circuito?, ¿qué mejora es la próxima y cuál puede esperar?\par
\textbf{Señal de progreso:} el equipo diferencia “posibilidades futuras” de “siguiente hito” y elige
un avance verificable.
\end{softcard}

\begin{card}
\labeltext{Cómo modelar el trabajo con IA}
Usar una estructura observable: \textbf{contexto + placa + componente + objetivo + restricciones +
formato de respuesta}. Pedir pinout o fuente técnica, contrastar la respuesta, revisar voltajes y
probar por partes. Si algo falla, describir el síntoma y cambiar una variable por vez. La evidencia
vale más que una respuesta convincente.
\end{card}

\begin{softcard}{AIEPRed}
\labeltext{Evitar}
Entregar una solución completa para copiar; reducir el taller a comandos o definiciones; energizar
sin revisión; presentar el PCB como punto de partida; prometer que toda idea será fabricada; o
confundir la exploración inicial con el desarrollo completo que ocurrirá después.
\end{softcard}

\newpage
% 6 ---------------------------------------------------------------------------
\kicker{Evaluación formativa y continuidad}
{\LARGE\bfseries\color{AIEPNavy}Cerrar con evidencia y una ruta de avance}\par
\vspace{4mm}

\begin{card}
\labeltext{Criterios de revisión de la ficha}
\begin{tabularx}{\textwidth}{@{}L{42mm}Y@{}}
\textbf{Coherencia} & El sensor elegido permite observar la variable asociada al problema.\\
\textbf{Funcionalidad} & La regla “cuando..., entonces...” describe una respuesta comprensible.\\
\textbf{Seguridad} & La prueba considera placa, voltaje, polaridad, resistencia y revisión previa.\\
\textbf{Verificación} & Existe evidencia inicial o un procedimiento de prueba definido con precisión.\\
\textbf{Autonomía} & El equipo registra qué aceptó, corrigió o descartó de la asistencia de IA.\\
\textbf{Continuidad} & El próximo hito puede comprobarse antes de la siguiente instancia.\\
\end{tabularx}
\end{card}

\begin{softcard}{AIEPGreen}
\labeltext{Escala rápida de avance}
\begin{tabularx}{\textwidth}{@{}L{28mm}Y@{}}
\textbf{Iniciado} & Hay problema e idea, pero la variable o el sensor aún no son claros.\\
\textbf{Encaminado} & Variable, sensor y regla son coherentes; falta prueba o evidencia.\\
\textbf{Verificable} & Existe evidencia inicial y el equipo puede explicar qué comprobó.\\
\textbf{Proyectado} & Además, el equipo prioriza una mejora y define su siguiente hito.\\
\end{tabularx}
\end{softcard}

\begin{card}
\labeltext{Ruta posterior al taller}
\begin{tabularx}{\textwidth}{@{}L{33mm}Y@{}}
\textbf{Mentoría 1 · 31 ago} & Validar el problema y comprender el contexto afectado.\\
\textbf{Mentoría 2 · 7 sep} & Afinar solución, recursos y factibilidad.\\
\textbf{Mentoría 3 · 21 sep} & Revisar prototipo, pruebas y evidencia.\\
\textbf{Mentoría 4 · 28 sep} & Construir relato, pitch y ensayo del equipo.\\
\textbf{Evento final · 5 oct} & Comunicar el proyecto y demostrar su avance.\\
\end{tabularx}
\end{card}

\begin{softcard}{AIEPCyan}
\labeltext{Sentido de las mentorías}
Las mentorías sirven para orientar, contrastar decisiones y desbloquear dificultades. El desarrollo
principal ocurre con el trabajo del equipo entre sesiones: investigar, construir, probar, documentar
y mejorar.
\end{softcard}

\begin{card}
\labeltext{Lista de cierre para quien facilita}
\begin{itemize}
  \item[$\square$] Cada equipo entregó una ficha identificable y legible.
  \item[$\square$] La variable y el sensor muestran una relación defendible.
  \item[$\square$] La regla inicial y el siguiente hito quedaron escritos.
  \item[$\square$] Cualquier prueba física fue revisada antes de energizar.
  \item[$\square$] Las fichas quedaron disponibles para la Mentoría 1.
\end{itemize}
\end{card}

\begin{softcard}{AIEPRed}
\centering
{\large\bfseries\color{AIEPNavy}GeoGreen demuestra hasta dónde puede crecer una idea. Hoy comienza la de ustedes.}
\end{softcard}

\end{document}
